\documentclass[a4paper]{article}
\usepackage[affil-it]{authblk}
\usepackage[english]{babel}
\usepackage[backend=bibtex,style=numeric]{biblatex}

\usepackage{geometry}
\geometry{margin=1.5cm, vmargin={0pt,1cm}}
\setlength{\topmargin}{-1cm}
\setlength{\paperheight}{29.7cm}
\setlength{\textheight}{25.3cm}

% useful packages.
\usepackage{amsfonts}
\usepackage{amsmath}
\usepackage{amssymb}
\usepackage{amsthm}     
\usepackage{enumerate} %enumerate environment
\usepackage{graphicx}
\usepackage{multicol}
\usepackage{fancyhdr}
\usepackage{layout}
\usepackage{ctex}      %中文环境
\usepackage{booktabs}  %三线表
\usepackage{verbatim}  %comment environment
\usepackage{float}     %强制表格图片位置
\usepackage{listings}  %insert code in latex
\usepackage{fontspec}  %font support
\usepackage{tikz}
\usepackage{makecell}
\usetikzlibrary{intersections}
\usepackage{arydshln}

\usepackage[colorlinks,linkcolor=blue]{hyperref}
\usepackage{xcolor}
\definecolor{mygreen}{rgb}{0,0.6,0}
\definecolor{lightgrey}{rgb}{0.95,0.95,0.95}
\definecolor{winered}{rgb}{0.5,0,0}
\definecolor{mymauve}{rgb}{0.58,0,0.82}
\lstset{
      backgroundcolor=\color{lightgrey},
      basicstyle             = \ttfamily,
      breakatwhitespace      = false,
      breaklines             = true,
      captionpos             = none,
      commentstyle           = \color{mygreen}\bfseries,
      extendedchars          = false,
      keepspaces             = true,
      keywordstyle           = \color{winered},         % keyword style
      language               = C++,                     % the language of code
      otherkeywords          = {string},
      numberstyle            = \tiny\color{mygray},
      rulecolor              = \color{black},
      showspaces             = false,
      showstringspaces       = false,
      showtabs               = false,
      stepnumber             = 1,
      stringstyle            = \color{mymauve},         % string literal style
      tabsize                = 2,
      title                  = \lstname,
      aboveskip              = 0pt,
      belowskip              = 0pt,
      columns                = flexible
}

% some common command
\newcommand{\dif}{\mathrm{d}}
\newcommand{\avg}[1]{\left\langle #1 \right\rangle}
\newcommand{\norm}[1]{\left\| #1 \right\|}
\newcommand{\difFrac}[2]{\frac{\dif #1}{\dif #2}}
\newcommand{\pdfFrac}[2]{\frac{\partial #1}{\partial #2}}
\newcommand{\OFL}{\mathrm{OFL}}
\newcommand{\UFL}{\mathrm{UFL}}
\newcommand{\fl}{\mathrm{fl}}
\newcommand{\op}{\odot}
\newcommand{\Eabs}{E_{\mathrm{abs}}}
\newcommand{\Erel}{E_{\mathrm{rel}}}
\addbibresource{citation.bib}

\begin{document}
% =================================================
\title{Quantum information and quantum computation homework 1}

\author{Zhouyue Qian 3220104074
  \thanks{Electronic address: \texttt{3220104074@edu.zju.cn}}}
\affil{(Information and Computational Science Class 2201), Zhejiang University }


\date{Due time: \today}

\maketitle

%\begin{abstract}
%    The abstract is not necessary for the theoretical homework, 
%    but for the programming project, 
%    you are encouraged to write one.      
%\end{abstract}





% ============================================
\section{Turing Machine 1}
I have built a Turing machine that can reserve the input string. It can reverse the input composed of "0"s and "1"s, similar to how "1001" would become "0110".
\subsection{\textbf{Design}}



\subsubsection{Symbols and States}

\begin{itemize}
    \item \textbf{External Alphabet \( A \)}: \(\{0, 1\}\)
    \item \textbf{Tape Alphabet \( S \)}: \(\{\textunderscore, 0, 1, *\}\), where \(*\) is an auxiliary marker.
    \item \textbf{State Set \( Q \)}: \(\{q_0, q_f, r, l, q_{Swap}, q_{Back}, q_{Skip}\}\), where:
    \begin{itemize}
        \item \( r \) indicates moving right.
        \item \( l \) indicates moving left.
        \item \( q_{Swap} \) indicates the swapping phase.
        \item \( q_{Back} \) indicates returning to the leftmost position after swapping.
        \item \( q_{Skip} \) indicates skipping consecutive blanks.
        \item \( q_f \) is the final state.
    \end{itemize}
\end{itemize}

\subsubsection{Transition Function}

\begin{enumerate}
    \item \textbf{Initial Phase (\( q_0 \))}: Mark the leftmost symbol and prepare to move right to find the rightmost symbol.
    \begin{align*}
        (q_0, 0) &\rightarrow (r, *, +1) \\
        (q_0, 1) &\rightarrow (r, *, +1) \\
        (q_0, \textunderscore) &\rightarrow (q_{Skip}, \textunderscore, +1)
    \end{align*}

    \item \textbf{Skipping Consecutive Blanks (\( q_{Skip} \))}: Skip consecutive blanks until a non-blank symbol is found.
    \begin{align*}
        (q_{Skip}, \textunderscore) &\rightarrow (q_{Skip}, \textunderscore, +1) \\
        (q_{Skip}, 0) &\rightarrow (r, 0, +1) \\
        (q_{Skip}, 1) &\rightarrow (r, 1, +1)
    \end{align*}

    \item \textbf{Moving Right (\( r \))}: Find the rightmost unprocessed symbol.
    \begin{align*}
        (r, 0) &\rightarrow (r, 0, +1) \\
        (r, 1) &\rightarrow (r, 1, +1) \\
        (r, \textunderscore) &\rightarrow (q_{Swap}, \textunderscore, -1)
    \end{align*}

    \item \textbf{Swapping Phase (\( q_{Swap} \))}: Move left to find the marker \(*\) and perform the symbol swap.
    \begin{align*}
        (q_{Swap}, 0) &\rightarrow (q_{Swap}, 0, -1) \\
        (q_{Swap}, 1) &\rightarrow (q_{Swap}, 1, -1) \\
        (q_{Swap}, *) &\rightarrow (q_{Back}, \textunderscore, +1)
    \end{align*}

    \item \textbf{Returning to the Start (\( q_{Back} \))}: Return to the leftmost position to process the next symbol.
    \begin{align*}
        (q_{Back}, 0) &\rightarrow (q_{Back}, 0, +1) \\
        (q_{Back}, 1) &\rightarrow (q_{Back}, 1, +1) \\
        (q_{Back}, *) &\rightarrow (r, *, +1)
    \end{align*}

    \item \textbf{Final State (\( q_f \))}: When all symbols are processed, stop.
    \begin{align*}
        (r, *) &\rightarrow (q_f, *, 0)
    \end{align*}
\end{enumerate}

\subsection{Advanced Analysis}
DThis Turing machine that can reverse an input binary string and handle consecutive blank spaces.
\subsubsection{Time Complexity}

\begin{itemize}
    \item \textbf{Initial Phase (\( q_0 \))}: \( O(1) \)
    \item \textbf{Moving Right (\( r \))}: \( O(n) \)
    \item \textbf{Swapping Phase (\( q_{Swap} \))}: \( O(n) \)
    \item \textbf{Returning to the Start (\( q_{Back} \))}: \( O(n) \)
    \item \textbf{Final State (\( q_f \))}: \( O(1) \)
\end{itemize}

\textbf{Total Time Complexity}: \( O(n^2) \), where \( n \) is the length of the input string.

\subsubsection{Space Complexity}

\begin{itemize}
    \item \textbf{Initial Phase (\( q_0 \))}: \( O(1) \)
    \item \textbf{Moving Right (\( r \))}: \( O(1) \)
    \item \textbf{Swapping Phase (\( q_{Swap} \))}: \( O(1) \)
    \item \textbf{Returning to the Start (\( q_{Back} \))}: \( O(1) \)
    \item \textbf{Final State (\( q_f \))}: \( O(1) \)
\end{itemize}

\textbf{Total Space Complexity}: \( O(n) \), where \( n \) is the length of the input string.

\section{Turing Machine 2}
This Turing machine can adds two numbers written in binary. 
\subsection{States}
\begin{itemize}
    \item $q_0$: Initial state. Reads the input and identifies two binary numbers separated by a delimiter "+". 
    \item $q_1$: Scanning the first binary number. 
    \item $q_2$: Carrying out the addition of corresponding bits from the two numbers.
    \item $q_3$: Writing the result and handling the carry.
    \item $q_4$: Final carry handling.
    \item $q_f$: Final state after completion.
\end{itemize}

\subsection{External Alphabet}
\begin{itemize}
    \item $0, 1$: Binary digits.
    \item $+$: Delimiter separating two numbers.
    \item $\_ $: Blank symbol.
\end{itemize}

\subsection{State Transition Function}

\begin{table}[h!]
\centering
\begin{tabular}{|c|c|c|c|c|}
    \hline
    \textbf{Current State} & \textbf{Read Symbol} & \textbf{Write Symbol} & \textbf{Move Direction} & \textbf{Next State} \\ \hline
    $q_0$ & 0 & 0 & R & $q_0$ \\ \hline
    $q_0$ & 1 & 1 & R & $q_0$ \\ \hline
    $q_0$ & + & + & R & $q_1$ \\ \hline
    $q_1$ & 0 & 0 & R & $q_1$ \\ \hline
    $q_1$ & 1 & 1 & R & $q_1$ \\ \hline
    $q_1$ & $\_$ & $\_$ & L & $q_2$ \\ \hline
    $q_2$ & 0 (no carry) & 0 & L & $q_3$ \\ \hline
    $q_2$ & 1 (no carry) & 1 & L & $q_3$ \\ \hline
    $q_2$ & 1 + 1 (no carry) & 0 & L & $q_3$ (carry) \\ \hline
    $q_3$ & 0 & 0 & L & $q_3$ \\ \hline
    $q_3$ & 1 & 1 & L & $q_3$ \\ \hline
    $q_3$ & + & + & L & $q_4$ \\ \hline
    $q_4$ & No carry & $\_$ & - & $q_f$ \\ \hline
    $q_4$ & Carry & 1 & L & $q_f$ \\ \hline
\end{tabular}
\caption{State Transition Table for Binary Addition}
\end{table} 

\subsection{Explanation of the Process}

1. In the initial state $q_0$, the Turing Machine scans the first binary number and stops when it encounters the "+" delimiter.
2. Once the "+" symbol is detected, the machine enters state $q_1$ and begins searching for the rightmost bit of the second binary number.
3. In state $q_2$, the machine adds the binary digits bit by bit, starting from the rightmost bit of both numbers.
4. The machine handles carry bits accordingly. If there is a carry after adding the bits, it enters state $q_3$ to continue the process.
5. In state $q_3$, the machine writes the result and moves left to handle the next bit or the carry.
6. If a final carry exists after all bits have been processed, the machine enters state $q_4$ and writes the carry before reaching the final state $q_f$.
7. The Turing Machine halts in state $q_f$ once all operations are complete. 

\subsection{Example: Adding 101 + 11}

For the input "101+11", the Turing Machine will proceed as follows:

\begin{itemize}
    \item Initial state $q_0$: Read "101", then encounter "+" and move to state $q_1$.
    \item State $q_1$: Start reading "11" from right to left.
    \item State $q_2$: Perform the addition:
    \begin{itemize}
        \item $1 + 1 = 10$: Write 0, carry 1.
        \item $0 + 1 + 1 = 10$: Write 0, carry 1.
        \item $1 + 0 + 1 = 10$: Write 0, carry 1.
    \end{itemize}
    \item State $q_4$: Write the final carry, resulting in "1000".
\end{itemize}

Thus, the binary addition of 101 and 11 yields the result 1000 (which equals 8 in decimal).

\section{Turing Machine proof}
\subsection{Basic Idea of the Simulation}

The basic idea of simulating a 2-tape Turing Machine \( M_2 \) with a single-tape Turing Machine \( M_1 \) involves encoding the two tapes of \( M_2 \) on a single tape of \( M_1 \). To do this, we need to:
\begin{itemize}
    \item Use an expanded alphabet for \( M_1 \) that can encode the symbols on both tapes of \( M_2 \), along with markers for the head positions.
    \item Simulate each step of \( M_2 \) by scanning the entire tape of \( M_1 \) to find the head positions, reading the symbols under each head, and then updating the tape according to the transition function of \( M_2 \).
\end{itemize}

\subsubsection{1-Tape Alphabet}
The alphabet of \( M_1 \) must be large enough to encode the symbols of both tapes of \( M_2 \). Let \( \Sigma \) be the alphabet of \( M_2 \), including the blank symbol \( \_\_ \). Then, each cell of \( M_1 \)'s tape must encode:
\begin{itemize}
    \item A symbol from \( \Sigma \) for tape 1 of \( M_2 \).
    \item A symbol from \( \Sigma \) for tape 2 of \( M_2 \).
    \item Two additional bits to mark the head positions on each of the two tapes.
\end{itemize}
Thus, the alphabet of \( M_1 \) is \( \Sigma^2 \times \{0, 1\}^2 \), which includes all possible combinations of symbols from both tapes and head markers.

\subsubsection{Initial Configuration}
Initially, the input on the first tape of \( M_2 \) is encoded on the single tape of \( M_1 \). The second tape of \( M_2 \) is initialized to be blank. The head positions of both tapes of \( M_2 \) are marked in the encoding on the single tape of \( M_1 \).

\subsection{Simulation Process}

The simulation proceeds by simulating each step of \( M_2 \) in two phases:
\begin{enumerate}
    \item A left-to-right scan of the tape to locate the head positions and read the symbols under both heads of \( M_2 \).
    \item A right-to-left scan to perform the actions dictated by the transition function of \( M_2 \) (write symbols and move the heads).
\end{enumerate}

\subsubsection{Phase 1: Left-to-Right Scan}
In this phase, \( M_1 \) performs a left-to-right scan of the entire tape to find the head positions of both tapes of \( M_2 \). During this scan, \( M_1 \) stores the symbols found under the heads of \( M_2 \) in its finite control. This takes \( O(s) \) steps, where \( s \) is the length of the portion of the tape that contains non-blank symbols.

\subsubsection{Phase 2: Right-to-Left Scan}
After reading the symbols under both heads in Phase 1, \( M_1 \) performs a right-to-left scan of the tape. During this scan, \( M_1 \) updates the symbols on the tape according to the transition function of \( M_2 \), moves the head markers, and performs any necessary write operations. This also takes \( O(s) \) steps.

Since each phase takes \( O(s) \) steps, the total time for simulating one step of \( M_2 \) is \( O(s) \).

\subsection{Time Complexity Analysis}

\subsubsection{Tape Length}
Let \( s \) denote the length of the portion of the tape used by \( M_2 \). In the worst case, the tape length grows with the number of steps \( T(n) \), so \( s \leq T(n) + 1 \).

\subsubsection{Total Time Complexity}
Each simulated step of \( M_2 \) requires \( O(s) \) steps by \( M_1 \). Since \( M_2 \) runs for \( T(n) \) steps, the total time required by \( M_1 \) to simulate all steps of \( M_2 \) is:
\[
O(s \cdot T(n)) = O(T(n) \cdot T(n)) = O(T(n)^2).
\]
Thus, the total time complexity of simulating a 2-tape Turing Machine with a single-tape Turing Machine is \( O(T(n)^2) \).  \hfill $\Box$

% ===============================================
%\section*{ \center{\normalsize {Acknowledgement}} }
%\printbibliography


\end{document}